% !TeX encoding = UTF-8
\documentclass[variant=docente,cover=institutional,mode=screen]{ua-engineering-v6-article}
% !TeX encoding = UTF-8
\UASetUniversity{Universidad Austral}
\UASetFaculty{Facultad de Ingeniería}
\UASetDepartment{Departamento de Computación}
\UASetCampus{Pilar}
\UASetCareer{Ingeniería Informática}
\UASetCommission{Comisión A}
\UASetProfessor{Equipo docente}
\UASetSemester{Segundo cuatrimestre 2026}
\UASetCourse{Sistemas Distribuidos}
\UASetDate{2026}


\UASetClassNumber{14}
\UASetClassTitle{Incidentes y cultura de confiabilidad}
\UASetDocKind{Guía docente}

\title{Clase 14 --- Guía docente}
\author{\UAProfessor}
\date{\UADate}

\begin{document}
\maketitle

\section{Objetivo pedagógico profundo}

La meta de esta clase es que el estudiante comprenda que la confiabilidad no se juega solo en arquitectura, observabilidad o resiliencia, sino también en la capacidad organizacional de responder y aprender de incidentes.

La idea que debe quedar instalada es:

\begin{center}
\textbf{la confiabilidad madura es una propiedad conjunta del sistema y de la organización que aprende de él}
\end{center}

\section{Resultados de aprendizaje esperados}

Al finalizar, el estudiante debería poder:

\begin{itemize}
\item definir incidente desde una perspectiva operativa;
\item describir una gestión razonable de incidentes;
\item explicar qué es un buen postmortem;
\item justificar por qué blameless no equivale a “nadie responde por nada”;
\item distinguir fix local de aprendizaje organizacional.
\end{itemize}

\section{Estructura sugerida para 90 minutos}

\subsection*{Bloque 1 (15 min): normalizar el incidente}
Abrir con una pregunta:
\begin{quote}
¿Qué revela un incidente: un error aislado o una propiedad del sistema y la organización?
\end{quote}

La discusión inicial debería llevar a abandonar la idea de que el incidente es solo una vergüenza técnica.

\paragraph{Objetivo didáctico}
Mover la conversación desde culpa puntual hacia aprendizaje sistémico.

\subsection*{Bloque 2 (20 min): gestión de incidentes}
Explicar fases:
\begin{itemize}
\item detección;
\item clasificación;
\item respuesta;
\item mitigación;
\item recuperación;
\item análisis posterior.
\end{itemize}

Introducir roles de coordinación.

\paragraph{Objetivo didáctico}
Mostrar que la respuesta al incidente también es una arquitectura organizacional.

\subsection*{Bloque 3 (20 min): postmortems}
Trabajar qué debe contener un buen postmortem y por qué el timeline es central.

\paragraph{Objetivo didáctico}
Que el alumno vea el postmortem como mecanismo de aprendizaje, no como burocracia retrospectiva.

\subsection*{Bloque 4 (15 min): blameless y causalidad}
Explicar por qué “culpa humana” como explicación final empobrece el análisis.

\paragraph{Objetivo didáctico}
Instalar una visión sistémica de causalidad distribuida y organizacional.

\subsection*{Bloque 5 (20 min): acciones y cultura}
Discutir qué acciones realmente cambian capacidad futura y cuáles son superficiales.

\paragraph{Objetivo didáctico}
Conectar postmortems con mejora real de arquitectura, observabilidad, procesos y cultura.

\section{Qué explicar y por qué}

\subsection*{1. Incidentes}
Porque es la unidad concreta donde se prueba la madurez operativa del equipo.

\subsection*{2. Gestión de incidentes}
Porque sin buena gestión, incluso sistemas razonables pueden degradarse muy mal.

\subsection*{3. Postmortems}
Porque son el principal mecanismo de aprendizaje explícito posterior al incidente.

\subsection*{4. Blameless}
Porque corrige una desviación cultural muy común y muy destructiva.

\subsection*{5. Aprendizaje organizacional}
Porque la clase debe mostrar que confiabilidad también es una capacidad colectiva.

\section{Preguntas disparadoras recomendadas}

\begin{itemize}
\item ¿Qué parte del incidente fue técnica y qué parte fue organizacional?
\item ¿Qué señal faltó para detectar antes?
\item ¿Qué decisión de mitigación funcionó realmente?
\item ¿Qué habría cambiado si el equipo hubiera tenido mejor runbook?
\item ¿Qué parte del incidente podría repetirse mañana si no cambia nada?
\end{itemize}

\section{Errores frecuentes del alumnado}

\begin{itemize}
\item creer que postmortem = encontrar culpable;
\item pensar que “blameless” significa ausencia de exigencia;
\item reducir el incidente al bug más visible;
\item proponer acciones vagas sin owner;
\item subestimar el valor del timeline.
\end{itemize}

\section{Ejercicio en vivo sugerido}

Proponer un incidente ficticio:
\begin{itemize}
\item alta latencia en DB;
\item retries en cascada;
\item API cae;
\item alertas tardías;
\item mitigación parcial tras 25 minutos.
\end{itemize}

Pedir al curso que reconstruya:
\begin{itemize}
\item impacto;
\item timeline;
\item factores contribuyentes;
\item acciones inmediatas;
\item acciones de mejora.
\end{itemize}

\paragraph{Objetivo}
Que el alumno practique el pasaje desde síntoma visible hacia aprendizaje estructurado.

\section{Momento crítico de la clase}

El punto más importante ocurre cuando el estudiante entiende que:

\begin{center}
\textbf{el valor real del incidente no está solo en restaurar servicio, sino en transformar la experiencia en capacidad organizacional futura}
\end{center}

Ese punto debe enfatizarse explícitamente.

\section{Criterio de éxito docente}

La clase salió bien si el estudiante puede:

\begin{itemize}
\item explicar qué distingue un buen postmortem de uno superficial;
\item construir una cadena causal razonable;
\item proponer acciones correctivas útiles y verificables;
\item justificar el rol de una cultura blameless madura;
\item conectar incidentes con confiabilidad y aprendizaje organizacional.
\end{itemize}

\section{Conexión con la clase siguiente}

Esta clase deja una transición natural hacia el cierre del curso, donde puede discutirse cómo todos los temas se integran en el diseño y defensa de un sistema distribuido completo.

\begin{uakeybox}
Pregunta puente: si ya entendemos mecanismos, arquitectura, observabilidad, resiliencia, objetivos e incidentes, ¿cómo defendemos integralmente un sistema distribuido ante un jurado técnico o comité de diseño?
\end{uakeybox}


\section{Plan de conducción ampliado}
La clase debe alternar explicación, predicción y contraste. Antes de revelar cada mecanismo, pedir al grupo que anticipe qué puede salir mal; después, volver al mismo escenario y preguntar qué propiedad mejora y qué costo aparece. Cada bloque conceptual debe cerrar con una decisión de diseño observable.
\subsection*{Secuencia recomendada}
\begin{enumerate}
\item Recuperar el prerequisito de la clase anterior.
\item Presentar un caso mínimo que falle y solicitar hipótesis.
\item Formalizar el concepto central de Incidentes y cultura de confiabilidad, distinguiendo seguridad, progreso y supuestos.
\item Resolver un ejemplo completo en pizarra, incluyendo una falla intermedia.
\item Comparar una alternativa de diseño y explicitar trade-offs.
\item Cerrar con una pregunta del Trabajo Final y una pregunta de defensa oral.
\end{enumerate}
\section{Diagnóstico de comprensión durante la clase}
No preguntar solamente ``¿se entendió?''. Usar preguntas discriminantes: ``¿qué cambia si el mensaje llega dos veces?'', ``¿qué propiedad sigue siendo cierta si cae un nodo?'', ``¿esto demuestra seguridad o progreso?'', ``¿qué observa un cliente externo?''. Si el estudiante responde solo con nombres de patrones, pedir una ejecución concreta.
\section{Criterios para corregir ejercicios}
Una respuesta excelente declara supuestos, construye una secuencia verificable, nombra la garantía con precisión, reconoce un costo y conecta la decisión con el dominio. Penalizar afirmaciones absolutas sin condiciones.
\section{Vinculación con el Trabajo Final Integrador}
Reservar entre 5 y 10 minutos para que cada grupo registre una decisión con: problema, alternativa elegida, alternativa descartada y razón. Esto crea trazabilidad entre las clases y las entregas acumulativas.
\section{Lectura docente previa}
Revisar la bibliografía indicada en los apuntes y preparar al menos un contraejemplo que no aparezca literalmente en las diapositivas. El objetivo es poder responder preguntas de frontera sin convertir la clase en una lectura del deck.

\section{Arquitectura didáctica v9 para 180 minutos}
\subsection*{Bloque 1 - desafío inicial (20 min)}
Presentar una situación donde aparezca detección tardía, coordinación confusa, causa simplificada y acciones sin owner. Pedir predicciones individuales antes de introducir vocabulario. El objetivo es hacer visibles los supuestos intuitivos y recuperar el prerequisito de la clase anterior.

\subsection*{Bloque 2 - construcción conceptual (35 min)}
Formalizar incident command, timeline, postmortems blameless, factores contribuyentes y aprendizaje organizacional. Separar seguridad, progreso y experiencia del usuario. Explicar no solo \emph{qué} hace cada mecanismo, sino \emph{por qué} existe y qué supuesto permite que funcione.

\subsection*{Bloque 3 - modelo y figura (35 min)}
Construir en pizarra una línea temporal, grafo, log, máquina de estados o tabla de decisiones. Recorrer una ejecución nominal y una adversarial. La representación debe quedar como referencia para los apuntes y el ejercicio principal.

\subsection*{Pausa (10 min)}
Recoger una pregunta anónima o un punto de confusión. Elegir una pregunta que permita distinguir síntoma, causa y propiedad.

\subsection*{Bloque 4 - caso trabajado con falla (40 min)}
Aplicar el mecanismo al caso conductor. Introducir una falla a mitad de ejecución y detenerse en cada transición: quién sabe qué, qué se persiste, qué garantía sigue vigente y qué observa el cliente.

\subsection*{Bloque 5 - taller de decisión (25 min)}
Los grupos aplican P-F-G-S-T a su Trabajo Final. Deben producir un ADR breve con alternativa elegida, alternativa descartada y evidencia pendiente. La evidencia de esta clase es timeline, impacto, factores contribuyentes y acciones verificables.

\subsection*{Bloque 6 - cierre y exit ticket (15 min)}
Cada estudiante responde: propiedad principal, supuesto decisivo, costo aceptado y condición bajo la cual cambiaría de diseño. Usar las respuestas para iniciar la clase siguiente.

\section{Qué explicar, por qué y cómo verificarlo}
\begin{itemize}
\item Explicar el problema antes del producto, porque reconocer la familia de problema es el resultado cognitivo central.
\item Explicar el invariante, porque permite evaluar configuraciones y no solo repetir un flujo nominal.
\item Explicar el límite, porque detección tardía, coordinación confusa, causa simplificada y acciones sin owner puede invalidar supuestos o reducir progreso.
\item Verificar comprensión mediante timeline, impacto, factores contribuyentes y acciones verificables, no mediante una definición aislada.
\end{itemize}

\section{Preguntas socráticas}
\begin{itemize}
\item ¿Qué sabe cada actor en este punto de la ejecución?
\item ¿Qué propiedad sigue siendo cierta si la respuesta nunca llega?
\item ¿Qué parte es safety y cuál es liveness?
\item ¿Qué alternativa más simple resolvería el problema si relajamos una garantía?
\item ¿Qué observación refutaría nuestra hipótesis de diseño?
\end{itemize}

\section{Errores previsibles y estrategia de corrección}
\begin{itemize}
\item Si el alumno responde con un nombre de tecnología, pedir actores, mensajes y estados.
\item Si confunde timeout con falla, construir dos ejecuciones indistinguibles.
\item Si promete todas las garantías, pedir comportamiento bajo partición o saturación.
\item Si mide solo rendimiento, preguntar cómo evidencia la propiedad semántica.
\item Si la solución es excesiva, pedir la alternativa mínima y el costo marginal de la garantía fuerte.
\end{itemize}

\section{Criterios de evaluación formativa}
Una respuesta excelente declara supuestos, recorre una ejecución, nombra con precisión la garantía, reconoce el costo y propone evidencia reproducible. Una respuesta suficiente identifica el mecanismo y el trade-off principal. Una respuesta insuficiente usa slogans, omite fallas o confunde producto con propiedad.

\section{Preparación docente y bibliografía}
Lectura previa recomendada: Google SRE Workbook, incident response y postmortems; literatura de safety science. Preparar un contraejemplo adicional y una variante del caso en la que cambie el costo del error; esto evita convertir la clase en una lectura de diapositivas.

\end{document}
